\chapter{Marco Teórico}

El desarrollo de un sistema informático para la gestión integral de establecimientos farmacéuticos demanda un sustento conceptual interdisciplinario que amalgama principios rigurosos de la ingeniería de software, la teoría de bases de datos relacionales, la logística de inventarios perecederos y la normatividad sanitaria aplicable a insumos para la salud. Este capítulo expone los fundamentos teóricos y metodológicos que respaldan la concepción, diseño y construcción del sistema \textit{Vitalis PharmaSys}.

\section{Arquitectura de Software y Patrones de Diseño Web}

En el ámbito de los sistemas de información contemporáneos, la adopción de patrones arquitectónicos desacoplados constituye una práctica indispensable para garantizar atributos de calidad tales como la mantenibilidad, escalabilidad, seguridad y testeabilidad del software \cite{pressman2014, sommerville2011}.

\subsection{El Patrón Arquitectónico Modelo-Vista-Controlador (MVC)}
El patrón Modelo-Vista-Controlador postula la separación de la aplicación en tres capas de responsabilidad claramente delimitadas:
\begin{itemize}
    \item \textbf{Modelo:} Encapsula la lógica de dominio, las entidades de negocio y las operaciones de acceso a datos. En el contexto de \textit{Vitalis PharmaSys}, los modelos representan abstracciones de medicamentos, lotes, recetas, ventas, compras y bitácoras operacionales.
    \item \textbf{Vista:} Corresponde a la capa de presentación que interactúa con el usuario final, encargada de renderizar la interfaz gráfica a través de plantillas HTML dinámicas generadas mediante el motor Blade de Laravel.
    \item \textbf{Controlador:} Actúa como intermediario orquestador; intercepta las solicitudes HTTP entrantes enviadas por el cliente, valida los datos de entrada a través de objetos de petición específicos (\textit{Form Requests}), invoca las operaciones de negocio pertinentes en los modelos y determina la vista o respuesta serializada que debe ser retornada.
\end{itemize}

\subsection{El Patrón Active Record y Mapeo Objeto-Relacional (ORM)}
El acceso a la persistencia se estructura a través del mapeador objeto-relacional Eloquent, el cual implementa el patrón de diseño \textit{Active Record} \cite{otwell2024laravel}. Bajo este enfoque, cada tabla del esquema relacional está asociada de manera unívoca a una clase de modelo, y cada fila de dicha tabla se representa como una instancia orientada a objetos. Este mecanismo abstrae la escritura de consultas SQL nativas, proveyendo constructores de consultas protegidos por parametrización obligatoria de sentencias, lo cual suprime las vulnerabilidades de inyección SQL y facilita la definición explícita de relaciones relacionales de cardinalidad uno a muchos ($1:N$) y muchos a muchos ($N:M$).

\section{Gestión de Bases de Datos Relacionales y Transacciones ACID}

La confiabilidad operativa de una farmacia depende críticamente de la consistencia matemática de sus inventarios y finanzas. Para ello, el repositorio integra el motor MariaDB bajo la conformación de tablas InnoDB.

\subsection{Normalización Relacional}
El diseño de la base de datos se estructuró conforme a las reglas de normalización formal hasta la Tercera Forma Normal (3FN) \cite{elmasri2016}:
\begin{enumerate}
    \item \textbf{Primera Forma Normal (1FN):} Se eliminan grupos repetitivos y atributos multivaluados, garantizando la atomicidad de cada columna (e.g., desglosando componentes de dosis y presentaciones en campos atómicos).
    \item \textbf{Segunda Forma Normal (2FN):} Estando en 1FN, todos los atributos no clave dependen funcionalmente de manera completa de la clave primaria, eliminando dependencias parciales en tablas con claves compuestas.
    \item \textbf{Tercera Forma Normal (3FN):} Estando en 2FN, se erradican las dependencias transitivas, asegurando que ningún atributo no clave dependa de otro atributo no clave. Por ejemplo, los datos del laboratorio o proveedor no se almacenan redundantemente en la tabla de medicamentos, sino que se vinculan mediante identificadores foráneos forzados por restricciones de integridad referencial (\textit{Foreign Keys}).
\end{enumerate}

\subsection{Propiedades Transaccionales ACID}
En operaciones concurrentes de dispensación y cobro, una venta involucra simultáneamente: la creación de la cabecera de venta, la inserción de partidas en el detalle, el descuento de existencias sobre lotes específicos, el registro del comprobante de receta médica y el asiento contable en bitácora. Para asegurar la integridad absoluta ante fallos imprevistos o concurrencia de terminales, estas operaciones se ejecutan bajo bloques transaccionales que satisfacen las propiedades ACID \cite{elmasri2016}:
\begin{itemize}
    \item \textbf{Atomicidad (Atomicity):} La transacción se ejecuta como una unidad indivisible; si cualquier consulta o validación falla, se efectúa una reversión total (\textit{rollback}), impidiendo estados intermedios inconsistentes (e.g., cobrar un producto sin descontar su stock).
    \item \textbf{Consistencia (Consistency):} La base de datos transita únicamente de un estado válido a otro, respetando todas las reglas de integridad de datos y restricciones de cardinalidad.
    \item \textbf{Aislamiento (Isolation):} Las operaciones simultáneas de dos o más cajeros sobre el mismo lote se gestionan de forma aislada mediante bloqueos a nivel de fila (\textit{row-level locking}), previniendo condiciones de carrera (\textit{race conditions}).
    \item \textbf{Durabilidad (Durability):} Una vez confirmada la transacción mediante \textit{commit}, las modificaciones de inventario y saldos persisten de manera inmutable incluso ante caídas súbitas del suministro eléctrico.
\end{itemize}

\section{Logística Farmacéutica: Estrategias FEFO y FIFO}

El control de inventario farmacéutico presenta complejidades singulares debido a la naturaleza perecedera de los medicamentos y la estricta vigencia terapéutica de los principios activos.

\subsection{El Principio FEFO (First Expired, First Out)}
A diferencia de los comercios generales donde prevalece la política FIFO (\textit{First In, First Out}), la administración farmacéutica exige primordialmente la aplicación de la regla FEFO (Primero en Expirar, Primero en Salir) \cite{cofepris2020}. Esta estrategia logística estipula que la prioridad de rotación y venta de un medicamento no está determinada por la fecha en que el producto ingresó a la bodega, sino por su fecha de caducidad certificada por el fabricante.

\begin{equation}
    \text{Prioridad de Despacho} = \arg\min_{l \in L_{m}} (\text{FechaCaducidad}_{l})
    \label{eq:fefo}
\end{equation}

Donde $L_{m}$ representa el conjunto de lotes activos con existencia disponible para el medicamento $m$. En \textit{Vitalis PharmaSys}, el motor de ventas consulta los lotes disponibles ordenados ascendentemente por fecha de vencimiento, asignando automáticamente las unidades a descontar del lote más próximo a caducar, minimizando sustancialmente las mermas por obsolescencia.

\section{Diseño de Interfaces de Usuario y Accesibilidad en Salud}

La interacción humana con sistemas clínicos demanda interfaces claras que minimicen la sobrecarga cognitiva del usuario en momentos de alta exigencia laboral.

\subsection{Filosofía de Diseño Clinical Precision UI}
El sistema adopta el paradigma de diseño \textit{Clinical Precision UI}, fundamentado en tres pilares de interacción:
\begin{itemize}
    \item \textbf{Espacio Full-Canvas:} Sustituye las barras laterales estáticas y pesadas por una cuadrícula fluida de 12 columnas que aprovecha el 100\% del ancho de la pantalla, incrementando el área visible para tablas de inventarios y detalles de facturación.
    \item \textbf{Lanzador de Aplicaciones Global (9 Puntos):} Menú flotante desplegable en cuadrícula $3 \times 3$ que provee acceso contextual instantáneo a los módulos neurálgicos (Dashboard, Inventario, Ventas, Recetas, Pacientes, Compras, Reportes, Usuarios y Ajustes) sin saturar la vista activa.
    \item \textbf{Paleta Cromática Clínica Controlada:} Uso de tonos petróleo (\#022B3A) para elementos institucionales, cerceta clínico (\#1F7A8C) para acciones interactivas primarias y cielo pálido (\#BFDBF7) para acentos informativos suaves, cumpliendo con las pautas de contraste de las directrices WCAG 2.1 nivel AA \cite{w3c2018wcag}.
\end{itemize}

\section{Marco Regulatorio Sanitario en México}

El funcionamiento de un sistema farmacéutico no puede concebirse al margen de las disposiciones emanadas por la Secretaría de Salud y la Comisión Federal para la Protección contra Riesgos Sanitarios (COFEPRIS):

\subsection{Trazabilidad de Lotes y Fechas de Caducidad (NOM-072-SSA1-2012)}
La Norma Oficial Mexicana NOM-072-SSA1-2012 establece los requerimientos obligatorios de etiquetado para medicamentos, estipulando la necesidad de identificar inequívocamente el número de lote de producción y la fecha de caducidad en el empaque primario y secundario \cite{nom072ssa1}. El sistema refleja esta exigencia en su modelo de datos al no permitir el ingreso de mercancía sin la vinculación estricta de dichos atributos.

\subsection{Dispensación y Retención de Recetas Médicas}
De acuerdo con los artículos aplicables de la Ley General de Salud y los lineamientos de COFEPRIS para antibióticos y fármacos controlados \cite{cofepris2020}:
\begin{enumerate}
    \item \textbf{Antibióticos (Fracción IV):} Requieren la retención y registro de la receta médica, la cual únicamente puede ser surtida durante el periodo de tratamiento prescrito por el facultativo.
    \item \textbf{Psicotrópicos y Estupefacientes (Fracciones I, II y III):} Exigen una fiscalización rigurosa que incluye el asiento ineludible del nombre del médico, número de cédula profesional, domicilio de prescripción, fecha de expedición y registro de la cantidad surtida. El módulo de dispensación de \textit{Vitalis PharmaSys} bloquea la confirmación de la venta si estos datos sanitarios no han sido debidamente capturados y validados.
\end{enumerate}
